\section{Methodik}
Da die vorliegende Arbeit im Rahmen der Lehrveranstaltung \enquote{Qualitative Forschung} durchgeführt wurde, lag die Wahl der Methodik auf qualitativen Forschungsmethoden. Die Wahl von qualitativen Forschungsmethoden ermöglicht es im Vergleich zu quantitativen Forschungsmethoden, tiefergehende und individuellere Einblicke in die Thematik zu bekommen. Hierbei konzentieren sich die gesammelten Daten auf subjektive Erfahrungen und Perspektiven der Interviewpartner.

Da es im Bereich dieses Themas bereits viele quantitative Studien gab und der Bereich der qualitativen Forschung weniger erforscht war, bot sich eine qualitative Herangehensweise sehr an, um das Thema aus einer ergänzenden, interpretativen Perspektive zu beleuchten.\indirekt{1}{mayring2000} \indirekt{39}{mayring2014}

Hierfür wurden mehrere strukturierte Interviews durchgeführt, da diese zwar eine systematische Herangehensweise bieten, aber dennoch eine offenere und flexiblerere Erfassung subjektiver Erfahrungen und Sichtweisen ermöglichen.\indirekt{257}{flick2022}

\subsection{Stichprobe \& Auswahlverfahren}
Die Auswahl der Interviewpartner erfolgte im Sinne des Convenience Sampling. Dabei wurde die Personen ausgewählt, die dem Forscher leicht zugänglich sind. Dabei ist diese Stichprobenwahl nicht repräsentativ und ermöglicht keine Verallgemeinerung der Ergebnisse auf die Gesamtpopulation. Aus diesem Grund ist das Convenience Sampling für explorative Studien, wie dieser gut geeignet.\indirekt{}{mayring2014}

Um aber bei dem Convenience Sampling und der begrenzten Stichprobengröße eine gewisse Varianz zu gewährleisten, wurden gezielt Interviewpartner aus unterschiedlichen Altersgruppen, Berufen und Geschlechtern gewählt. Es wurde auch darauf geschaut, dass sich die Interviewpartner bezüglich der Online-Affinität und des Konsumverhaltens unterscheiden. Ziel war es trotz der etwas begrenzten Auswahlmöglichkeit und einer wenig repräsentativen Stichprobe, dennoch eine gewisse Breite zu erreichen. Es ist anzumerken, dass aufgrund der limitierten Stichprobengröße die Übertragbarkeit der Ergebnisse auf die Gesamtpopulation somit eingeschränkt ist.

Insgesamt wurden im Rahmen dieser Arbeit drei Interviews durchgeführt.

\subsection{Datengenerierung}
Die Durchschnittsdauer der Interviews betrug etwa 10 Minuten. Alle Interviews wurde mit expliziter Zustimmung der Interviewpartner aufgenommen und anschließend transkribiert, um eine akkurate und detaillierte Analyse zu ermöglichen. Die Interviews wurden in Präsenz, in einer ruhigen und vertrauten Umgebung durchgeführt, um eine offene und ehrliche Kommunikation zu fördern.\indirekt{}{mayring2014}

Vor Beginn des Interviews wurde neben der Zustimmung zur Aufzeichung des Interviews auch über das Ziel und den Ablauf des Interviews informiert. Dies sollte den Interviewpartnern ermöglichen, sich selbst noch kurz auf das Interview vorzubereiten. Außerdem sollte so die Möglichkeit gegeben werden, die Einwilligung zum Interview wieder zurückzuziehen.\indirekt{}{bfdidatenschutz}

\subsection{Datenauswertung}
Die Auswertung der gesammelten Daten der Interviews erfolgte mittels der qualitativen Inhaltsanalyse nach Mayring. Dieses Verfahren ermöglicht ein schrittweises Vorgehen für die systematische Analyse und Interpretation von Texten. Dafür werden Kategorien gebildet, die auf den Inhalten der Texte basieren. Damit sollen die zentralen Themen und Muster der Texte diesen Kategorien zugordnet werden.\indirekt{1}{mayring2000}
Aufgrund der vorherigen Behandlung dieses Analyseverfahrens in der Lehrveranstaltung, wurde die qualitative Inhaltsanalyse nach Mayring als Methode für die Datenauswertung gewählt.

Das Ziel der Datenanalyse lag darin zentrale Muster in der Entwicklung des Kaufverhaltens in der Corona-Pandemie zu identifizieren. Für die Auswertung orientierte man sich dabei an einer deduktiv-induktiven Vorgehensweise.

\subsubsection{Vorbereitung des Materials}
Nun wurden die Interviews im vollen Umfang zunächst einmal mit dem Einsatz einer Speech-To-Text KI transkribiert. Das Ergebnis wurde dann noch mehrmals mit den Audioaufnahmen abgeglichen, um so sicherzustellen, dass das Transkript akkurat war. Im Laufe dieser Einarbeitung kam es dann noch zu einer Kennzeichnung von potentiell relevanten Textstellen, die als Kodiereinheit agierten sollen. 

\subsubsection{Entwicklung eines Kategoriensystems}
Um später mit der qualitativen Inhaltsanalyse beginnen zu können mussten Kategorien verfasst werden. Dafür orientierte man sich zunächst an der Forschungsfrage und überlegte sich erste Oberkategorien wie beispielsweise Bequemlichkeit, Atmosphäre oder Kosten. Nach der Betrachtung der Fforschungsfrage wurden dann die Transkripte und der Fragenkatalog noch einmal durchgearbeitet um gezielt Kategorien finden zu können.

Um die Erarbeitung von Kategorien besserer Qualität zu fördern wurde für eine Überarbeitung der Kategorien ein Large Language Model (LLM) als analytisches Hilfsmittel verwendet. Dieses KI-System konnte dann auf Basis der vorläufigen Kategorien sowie den Transkripten, dem Forschungskatalog und der Forschungsfrage einmal selbst eine Überarbeitung der Kategorien vornehmen. Eine Arbeit mit einem LLM war bereits sehr geläufig und wurde im Rahmen der Lehrveranstalung auch als Hilfsmittel empfohlen.

Nach Sichtung der überarbeiteten Kategorien wurden diese durch den Forschenden noch etwas weiter ausgearbeitet. Das LLM wurde also ausschließlich als Unterstützung bei der Erarbeitung verwendet und nicht als Ersatz für die interpretative Analyse. Aus den überarbeiteten Daten durch das LLM und den Forschenden entstand letztendlich ein finaler Kodierleitfaden mit Kategorien, Definitionen, Kodierregeln und Ankerbeispielen.

\subsubsection{Kodierungsprozess}
Im folgenden Schritt wurden die Transkripte dann systematisch kodiert wobei man relevanten Textpassagen des Interviews dann den Kategorien zuordnete. Es kam während des Kodierungsprozesses allerdings auch vor, dass bemerkt wurde, dass Kategorien eventuell nicht konkret genug oder etwas zu ungenau waren. In diesem Fall wurden die Kategorien dann noch geringfügig angepasst.
